\documentclass[11pt]{article}
\usepackage[margin=2.2cm]{geometry}
\usepackage{booktabs}
\usepackage{amsmath}
\usepackage{enumitem}
\usepackage[T1]{fontenc}
\usepackage{lmodern}

\setlength{\parindent}{0pt}
\setlength{\parskip}{6pt}
\newcommand{\code}[1]{\texttt{#1}}
\newcommand{\answerbox}[1]{\vspace{2pt}\fbox{\parbox[c][#1][c]{\dimexpr\linewidth-2\fboxsep-2\fboxrule}{\ }}\vspace{4pt}}

\begin{document}

\begin{center}
{\Large\bfseries EE\,362 --- Pre-Lab 1}\\[2pt]
{\large Design Flow and the Transparent Latch}\\[6pt]
\begin{tabular}{@{}p{6.5cm}p{4cm}p{4cm}@{}}
Name: \hrulefill & Section: Mon / Tue & Date: \hrulefill
\end{tabular}
\end{center}

\vspace{2pt}\hrule\vspace{4pt}

{\small\textbf{Due at the start of the lab session.} Five points. About twenty
minutes of work, on paper, before you open Quartus. The purpose is not to test
you --- it is to make sure you arrive with a design already in your head, so the
two hours at the bench are spent building instead of thinking.}

\vspace{4pt}\hrule

\section*{Question 1 --- Paper first}

Complete the transition table for a 1-bit transparent latch with an asynchronous
clear. Use \code{--} for \emph{don't care} and \textbf{Q} to mean \emph{hold the
previous value}.

\begin{center}
\begin{tabular}{clll@{\hskip 8pt}||@{\hskip 8pt}l@{\hskip 20pt}l}
\toprule
\# & \code{Clr} & \code{En} & \code{D} & $Q^{+}$ & note \\
\midrule
0 & 1 & -- & -- & \rule{2.2cm}{0.4pt} & \rule{4.5cm}{0.4pt} \\[4pt]
1 & 0 & 0 & -- & \rule{2.2cm}{0.4pt} & \rule{4.5cm}{0.4pt} \\[4pt]
2 & 0 & 1 & 0 & \rule{2.2cm}{0.4pt} & \rule{4.5cm}{0.4pt} \\[4pt]
3 & 0 & 1 & 1 & \rule{2.2cm}{0.4pt} & \rule{4.5cm}{0.4pt} \\
\bottomrule
\end{tabular}
\end{center}

Now write the characteristic equation for $Q^{+}$ in terms of
\code{Clr}, \code{En}, \code{D} and \code{Q}:

\answerbox{2.0cm}

\textbf{One sentence:} which term in your equation is the one that makes this
circuit \emph{memory} rather than combinational logic, and why?

\answerbox{1.6cm}

\section*{Question 2 --- Interface}

Write the VHDL \code{entity} declaration for \code{TransLatch} --- the interface
only, no architecture. Get the port names, directions and types right.

\answerbox{4.2cm}

Now the same for \code{TransLatch8}. \textbf{Before you write it:} which ports
change width, and which stay one bit wide?

\answerbox{4.2cm}

\section*{Question 3 --- Test plan}

List the minimum set of input vectors that would convince \emph{you} that a
\code{TransLatch} implementation is correct. For each one, say what it proves.

\begin{center}
\begin{tabular}{@{}c@{\hskip 8pt}p{5.2cm}@{\hskip 8pt}p{7.6cm}@{}}
\toprule
\# & Inputs (\code{Clr}, \code{En}, \code{D}) & What this vector proves \\
\midrule
1 & \rule{5cm}{0.4pt} & \rule{7.4cm}{0.4pt} \\[8pt]
2 & \rule{5cm}{0.4pt} & \rule{7.4cm}{0.4pt} \\[8pt]
3 & \rule{5cm}{0.4pt} & \rule{7.4cm}{0.4pt} \\[8pt]
4 & \rule{5cm}{0.4pt} & \rule{7.4cm}{0.4pt} \\[8pt]
5 & \rule{5cm}{0.4pt} & \rule{7.4cm}{0.4pt} \\[8pt]
6 & \rule{5cm}{0.4pt} & \rule{7.4cm}{0.4pt} \\
\bottomrule
\end{tabular}
\end{center}

\textbf{Then answer this, in one or two sentences:} why is your set
\emph{sufficient}? That is, what would a wrong implementation have to do in order
to pass all of your vectors and still be wrong?

\answerbox{2.2cm}

\vfill
{\small\emph{Question 3 is the one that matters.} It is the same question as
Section 5 of your lab report, and it is the difference between testing a design
and merely running it.}

\end{document}
